% Standalone problem document, generated from the adjacent README.md.
% Compile with XeLaTeX. Mathematical content is maintained in Markdown.
\documentclass[11pt,a4paper]{article}
\usepackage[fontsize=10.5pt]{scrextend}
\usepackage[margin=22mm,headheight=15pt,headsep=9mm,footskip=11mm]{geometry}
\usepackage{fontspec}
\setmainfont{texgyrepagella}[Extension=.otf,UprightFont=*-regular,BoldFont=*-bold,ItalicFont=*-italic,BoldItalicFont=*-bolditalic]
\setsansfont{texgyreheros}[Extension=.otf,UprightFont=*-regular,BoldFont=*-bold,ItalicFont=*-italic,BoldItalicFont=*-bolditalic]
\setmonofont{texgyrecursor}[Extension=.otf,UprightFont=*-regular,BoldFont=*-bold,ItalicFont=*-italic,BoldItalicFont=*-bolditalic,Scale=MatchLowercase]
\usepackage{amsmath,amssymb}
\usepackage{unicode-math}
\setmathfont{texgyrepagella-math.otf}
\usepackage{xcolor,microtype,enumitem,needspace,fancyhdr}
\usepackage{bookmark}
\definecolor{ink}{HTML}{17344B}
\definecolor{accent}{HTML}{197D80}
\definecolor{muted}{HTML}{596773}
\hypersetup{colorlinks=true,linkcolor=accent,urlcolor=accent,pdftitle={SP-03 - The
Euclidean distance degree of the real symplectic
group},pdfauthor={Open Problems in Numerical Linear Algebra}}
\urlstyle{same}
\setlength{\parindent}{0pt}
\setlength{\parskip}{5pt plus 1pt minus 1pt}
\linespread{1.04}
\setlength{\emergencystretch}{3em}
\widowpenalty=10000
\clubpenalty=10000
\displaywidowpenalty=10000
\allowdisplaybreaks[1]
\setlist{nosep,leftmargin=1.5em}
\providecommand{\tightlist}{\setlength{\itemsep}{0pt}\setlength{\parskip}{0pt}}
\providecommand{\pandocbounded}[1]{#1}
\pagestyle{fancy}
\fancyhf{}
\fancyhead[L]{\sffamily\footnotesize\color{muted}OPEN PROBLEMS IN NUMERICAL LINEAR ALGEBRA}
\fancyhead[R]{\sffamily\bfseries\footnotesize\color{accent}SP-03}
\fancyfoot[L]{\parbox[b]{0.9\textwidth}{\sffamily\fontsize{8}{10}\selectfont\color{muted}Editorial ratings; a bounded search cannot certify that a problem remains unsolved.\\Literature check: 2026-09-10}}
\fancyfoot[R]{\sffamily\footnotesize\color{muted}\thepage}
\renewcommand{\headrulewidth}{0.3pt}
\renewcommand{\headrule}{\hbox to\headwidth{\color{accent}\leaders\hrule height \headrulewidth\hfill}}
\makeatletter
\renewcommand{\section}{\Needspace{5\baselineskip}\@startsection{section}{1}{0pt}{12pt plus 2pt minus 2pt}{4pt}{\sffamily\bfseries\large\color{ink}}}
\renewcommand{\subsection}{\Needspace{4\baselineskip}\@startsection{subsection}{2}{0pt}{9pt plus 1pt minus 1pt}{3pt}{\sffamily\bfseries\normalsize\color{ink}}}
\makeatother
\setcounter{secnumdepth}{0}
\begin{document}
{\sffamily\small\color{muted}Eigenvalues and inverse problems\par}
\vspace{3pt}
{\raggedright\hyphenpenalty=10000\exhyphenpenalty=10000\sffamily\bfseries\fontsize{21}{24}\selectfont\color{ink}The
Euclidean distance degree of the real symplectic group\par}
\vspace{7pt}
{\sffamily\small\textbf{Difficulty:} challenging\quad\textbf{Importance:} interesting
to specialist\par}
{\sffamily\small\bfseries\color{accent}Status: Open\par}
\vspace{4pt}
{\color{accent}\hrule height 0.8pt}
\vspace{6pt}
\textbf{Rating rationale:} Challenging reflects an all-ranks algebraic
critical-point count with only low-rank computations; specialist impact
concerns the algebraic complexity of symplectic matrix nearness.

For each \(m\ge1\), set \[
J=\begin{pmatrix}0&I_m\\-I_m&0\end{pmatrix},\qquad
\operatorname{Sp}_{2m}(\mathbb C)=
\{X\in\mathbb C^{2m\times2m}:X^TJX=J\}.
\] For a generic data matrix \(U\in\mathbb C^{2m\times2m}\), let \(D_m\)
be the number of complex critical points on this variety of \[
f_U(X)=\operatorname{tr}\bigl((X-U)^T(X-U)\bigr).
\] This is the complexification of squared real Frobenius distance; the
transpose is not a conjugate transpose. Equivalently, count the
solutions \(X\) of \(X^TJX=J\) and \[
\operatorname{tr}\bigl((U-X)^TXH\bigr)=0
\quad\text{for every }H\in\mathbb C^{2m\times2m}
\text{ satisfying }H^TJ+JH=0.
\] ``Generic'' means outside a proper Zariski-closed set where the
finite critical-point count can change. Is \[
D_m=2^{m^2}+2^{\,2m-1}
\qquad\text{for every }m\ge1?
\]

\subsection{Relevance}\label{relevance}

\(D_m\) measures the algebraic complexity of finding the nearest real
symplectic matrix. This is a structured matrix nearness problem,
relevant when numerical approximations should preserve a symplectic
form. The standard Frobenius inner product and the displayed \(J\) are
part of the problem.

\subsection{References}\label{references}

\begin{itemize}
\tightlist
\item
  J. A. Baaijens and J. Draisma, \emph{Euclidean distance degrees of
  real algebraic groups}, Linear Algebra and its Applications 467
  (2015), 174--187, §5, p.187: proposed formula immediately before
  Problem 5.1; §2 for critical equations
  (\href{https://pure.tue.nl/ws/files/3846347/391917266748824.pdf}{version
  of record};
  \href{https://doi.org/10.1016/j.laa.2014.11.012}{journal}).
\item
  Z. Lai, L.-H. Lim, and K. Ye, \emph{Euclidean Distance Degree in
  Manifold Optimization}, SIAM Journal on Optimization 35 (2025),
  2402--2422, related ED-degree results for flag, Grassmann, and Stiefel
  models (\href{https://doi.org/10.1137/25M1735032}{journal}).
\end{itemize}

\subsection{Status check --- 2026-09-08}\label{status-check-2026-09-08}

The source computed \(D_1=4\), \(D_2=24\), and \(D_3=544\), and
explicitly proposed the displayed pattern. It did not assert a proof.
Searches for ``symplectic group'', ``Euclidean distance degree'', ``ED
degree'', ``544'', and 2024--2026 found no proof, counterexample, or
additional general formula. The 2025 related article's abstract treats
different manifold families. No recent explicit reaffirmation of this
exact formula was located.

\subsection{Audit update --- 2026-09-10}\label{audit-update-2026-09-10}

Rechecked
\href{https://pure.tue.nl/ws/files/3846347/391917266748824.pdf}{Baaijens--Draisma},
§5: the general expression is proposed from small-rank computations,
without a proof. Searches for symplectic Euclidean distance degrees and
subsequent matrix-manifold ED-degree work found no derivation of this
formula. Importance is narrowed to specialist because the requested
output is a specific algebraic degree, rather than a general symplectic
projection algorithm.
\end{document}
